\documentclass[12pt,a4paper]{article}

% -------------------------------------------------
% Sprache / Kodierung
% -------------------------------------------------
\usepackage[T1]{fontenc}
\usepackage[utf8]{inputenc}
\usepackage[ngerman]{babel}
\usepackage{lmodern}
\usepackage{microtype}

% -------------------------------------------------
% Mathematik
% -------------------------------------------------
\usepackage{amsmath,amssymb,amsthm,mathtools}

% -------------------------------------------------
% Layout / Grafik / Farben
% -------------------------------------------------
\usepackage[a4paper,margin=2.3cm]{geometry}
\usepackage{booktabs}
\usepackage{xcolor}
\usepackage[hidelinks]{hyperref}
\usepackage[most]{tcolorbox}
\usepackage{enumitem}
\usepackage{array}

% -------------------------------------------------
% Farben
% -------------------------------------------------
\definecolor{lückenrot}{RGB}{180,40,40}
\definecolor{bewiesengrün}{RGB}{30,100,50}
\definecolor{warnorange}{RGB}{200,100,10}
\definecolor{hintergrundblau}{RGB}{235,242,250}
\definecolor{hintergrundgelb}{RGB}{255,252,225}
\definecolor{fehlerrot}{RGB}{200,0,0}
\definecolor{konjviolett}{RGB}{80,0,130}

% -------------------------------------------------
% tcolorbox-Stile
% -------------------------------------------------
\tcbset{
  lücke/.style={colback=red!8, colframe=lückenrot, fonttitle=\bfseries,
    title={Offene Lücke}},
  kernidee/.style={colback=hintergrundblau, colframe=blue!60!black,
    fonttitle=\bfseries, title={Kernidee}},
  ergebnis/.style={colback=green!7, colframe=bewiesengrün,
    fonttitle=\bfseries, title={Bewiesenes Ergebnis}},
  warnung/.style={colback=orange!10, colframe=warnorange,
    fonttitle=\bfseries, title={Achtung: Heuristik\,/\,Konjektur}},
  korrektur/.style={colback=red!4, colframe=fehlerrot, fonttitle=\bfseries,
    title={Korrektur eines Vorgängerdokuments}},
  konj/.style={colback=violet!5, colframe=konjviolett, fonttitle=\bfseries,
    title={Konjektur (unbewiesen)}},
}

% -------------------------------------------------
% Theorem-Umgebungen
% -------------------------------------------------
\newtheorem{definition}{Definition}[section]
\newtheorem{proposition}[definition]{Proposition}
\newtheorem{theorem}[definition]{Theorem}
\newtheorem{lemma}[definition]{Lemma}
\newtheorem{korollar}[definition]{Korollar}
\newtheorem{bemerkung}[definition]{Bemerkung}
\newtheorem{hypothese}[definition]{Hypothese}
\newtheorem{satz}[definition]{Satz}

\newenvironment{beweis}{\begin{proof}[Beweis]}{\end{proof}}

% -------------------------------------------------
% Makros
% -------------------------------------------------
\newcommand{\vii}{\nu_2}
\newcommand{\U}{\mathcal{U}}
\newcommand{\N}{\mathbb{N}}
\newcommand{\Z}{\mathbb{Z}}
\newcommand{\R}{\mathbb{R}}
\newcommand{\Ztwo}{\mathbb{Z}_2}
\newcommand{\abs}[1]{\lvert#1\rvert}

% -------------------------------------------------
% Dokument
% -------------------------------------------------
\title{%
  Collatz-Beweisversuch: Synthese und Abschlussanalyse\\[0.4em]
  \large Kombinierter Abstieg, Zyklenfreiheit und ehrliche Bestandsaufnahme%
}
\author{Thomas Hoffbauer}
\date{\today}

\begin{document}
\maketitle

\begin{abstract}
Dieses Dokument synthetisiert die Ergebnisse aus fünf Vorgängertexten
(\textit{eabc\_spektrum}, \textit{oktonionen\_beweis},
\textit{dc\_morley\_walter}, \textit{c\_ketten\_2adisch},
\textit{lean\_dc}) zu einem Gesamtbild.

Tatsächlich rigoros bewiesen sind: (i)~B-Typ kann nicht auf B-Typ folgen,
(ii)~C-Ketten der Länge~$k$ erfordern Startwerte in einer Residuenklasse
der Dichte $\leq 1/2^k$, (iii)~eine kompakte Iterationsformel für
C-Kettenglieder, (iv)~Normdeszenz für EA-Schritte im $E_8$-Gitter für
hinreichend große Normen.

Als neues Ergebnis wird hier gezeigt: Ein vollständiger B\textbullet C\textbullet A-Block
mit~$\vii = 2$ beim A-Schritt wächst um Faktor $\approx 27/16$; Kontraktion
tritt genau dann ein, wenn $\vii \geq 3$ beim A-Schritt gilt.
Für das Zyklenausschluss-Argument wird die Tao-Identität mit der DC-Struktur
verbunden: Für einen Zyklus mit B\textbullet C\textbullet A-Blöcken gilt
$2^m < 3^l$ (nötig wäre $2^m > 3^l$) --- ein direkter Widerspruch,
der Zyklen dieser Struktur ausschließt.

Offen bleibt die Lücke zwischen dem 2-adischen Argument (keine unendlichen
C-Ketten in $\Ztwo$) und der universellen Aussage für alle $n \in \N$.
\end{abstract}

\tableofcontents

\bigskip
\begin{tcolorbox}[warnung, title={Zur Einordnung dieses Dokuments}]
Die Collatz-Vermutung ist unbewiesen. Dieses Dokument führt den Stand
des Beweisversuchs so weit wie möglich voran, trennt aber konsequent
zwischen \emph{bewiesenen Aussagen}, \emph{Konjekturen} und
\emph{offenen Lücken}. Ein Beweis der Collatz-Vermutung liegt hier nicht vor.
\end{tcolorbox}

% ====================================================
\section{Bewiesene Bausteine: Übersicht}
\label{sec:bausteine}
% ====================================================

Die folgende Tabelle fasst alle rigoros bewiesenen Ergebnisse aus
den fünf Vorgängerdokumenten zusammen.

\begin{center}
\renewcommand{\arraystretch}{1.45}
\begin{tabular}{>{\raggedright}p{5.8cm} >{\centering}p{2.8cm} >{\raggedright\arraybackslash}p{4.5cm}}
\toprule
\textbf{Aussage} & \textbf{Status} & \textbf{Quelle / Methode} \\
\midrule
$\vii(3n+1) = 1$ für $n \equiv 7,11 \pmod{12}$
  & \textcolor{bewiesengrün}{\bfseries Bewiesen}
  & EABC-Spektrum, Kongruenzrechnung \\[2pt]
$\vii(3n+1) \geq 2$ für $n \equiv 1,5 \pmod{12}$
  & \textcolor{bewiesengrün}{\bfseries Bewiesen}
  & EABC-Spektrum, Kongruenzrechnung \\[2pt]
B-Typ folgt nicht auf B-Typ
  & \textcolor{bewiesengrün}{\bfseries Bewiesen}
  & DC-Morley-Walter, Prop.\,2.1 \\[2pt]
BC-Ketten haben Form $\mathrm{B}^{\leq 1}\cdot\mathrm{C}^k$
  & \textcolor{bewiesengrün}{\bfseries Bewiesen}
  & DC-Morley-Walter, Kor.\,2.1 \\[2pt]
C-Ketten Länge~$k$: Dichte $\leq 1/2^k$
  & \textcolor{bewiesengrün}{\bfseries Bewiesen}
  & DC-Morley-Walter, Induktion \\[2pt]
Formel $n_j = 3^j(n_0+1)/2^j - 1$
  & \textcolor{bewiesengrün}{\bfseries Bewiesen}
  & C-Ketten-2-adisch, Prop.\,2.1 \\[2pt]
C-Kette Länge $k$ erfordert $\vii(n_0+1) \geq k$
  & \textcolor{bewiesengrün}{\bfseries Bewiesen}
  & C-Ketten-2-adisch, Prop.\,2.2 \\[2pt]
Unendliche C-Kette $\Rightarrow n_0 = -1 \in \Ztwo$
  & \textcolor{bewiesengrün}{\bfseries Bewiesen}
  & 2-adischer Limes \\[2pt]
$E_8$-Normformel: $N(3x+1) = 9N(x)+6\operatorname{Re}(x)+1$
  & \textcolor{bewiesengrün}{\bfseries Bewiesen}
  & Oktonionen, Prop.\,3.1 \\[2pt]
Normabstieg bei EA-Schritt ($N \gg 1$)
  & \textcolor{bewiesengrün}{\bfseries Bewiesen}
  & Oktonionen-Normabschätzung \\[2pt]
Tao-Identität: $\U^l(n)\cdot 2^m = 3^l\cdot n + G_l$
  & \textcolor{warnorange}{\bfseries extern}
  & johnjanik/syracuse \\[2pt]
Keine positiven Zyklen (bedingt)
  & \textcolor{warnorange}{\bfseries bedingt}
  & ericmerle, Baker-Axiom \\[2pt]
DC (gleichmäßige Schranke $m_0$)
  & \textcolor{lückenrot}{\bfseries offen}
  & --- \\[2pt]
Collatz-Vermutung
  & \textcolor{lückenrot}{\bfseries offen}
  & --- \\
\bottomrule
\end{tabular}
\end{center}

% ====================================================
\section{EABC-Grundstruktur: Präzisierung}
\label{sec:eabc}
% ====================================================

\subsection{Korrekte mod-24-Aussagen}

\begin{tcolorbox}[korrektur]
Das Dokument \texttt{collatz\_lean\_dc.tex} enthält eine missverständliche
Aussage: Die Lean\,4-Verifikation
\[
  \texttt{forall r : ZMod 12,\; r = 11 -> (3*r+1)/2 = 5}
\]
prüft lediglich, dass für den \emph{Repräsentanten} $r = 11$ in \texttt{ZMod\,12}
der Ausdruck $34/2 = 17 \equiv 5 \pmod{12}$ gilt.
Das bedeutet \emph{nicht}, dass $\U(n) \equiv 5 \pmod{12}$ für alle
$n \equiv 11 \pmod{12}$. Gegenbeispiel: $n = 23$, $\U(23) = 35 \equiv 11 \pmod{12}$.

Die korrekte Aussage benötigt Modulus~$24$, nicht~$12$.
\end{tcolorbox}

\begin{proposition}[Übergangsregeln modulo 24]
\label{prop:mod24}
Für ungerades $n \in \N$ gilt unter der odd-to-odd-Abbildung $\U(n) = (3n+1)/2^{\vii(3n+1)}$:
\begin{align*}
  n \equiv 7  \pmod{24} \;&\Rightarrow\; \U(n) \equiv 11 \pmod{12},\quad (\text{B} \to \text{C})\\
  n \equiv 19 \pmod{24} \;&\Rightarrow\; \U(n) \equiv 5  \pmod{12},\quad (\text{B} \to \text{A})\\
  n \equiv 11 \pmod{24} \;&\Rightarrow\; \U(n) \equiv 5  \pmod{12},\quad (\text{C} \to \text{A})\\
  n \equiv 23 \pmod{24} \;&\Rightarrow\; \U(n) \equiv 11 \pmod{12},\quad (\text{C} \to \text{C})
\end{align*}
\end{proposition}

\begin{beweis}
Sei $n = 24k + r$ für $r \in \{7,19,11,23\}$.
\textit{Fall $r=7$}: $n_1 = (3(24k+7)+1)/2 = (72k+22)/2 = 36k+11$.
$36k+11 \equiv 11 \pmod{12}$ (da $36k \equiv 0 \pmod{12}$). \\
\textit{Fall $r=19$}: $n_1 = (72k+58)/2 = 36k+29$. $36k+29 \equiv 5 \pmod{12}$. \\
\textit{Fall $r=11$}: $n_1 = (72k+34)/2 = 36k+17$. $36k+17 \equiv 5 \pmod{12}$. \\
\textit{Fall $r=23$}: $n_1 = (72k+70)/2 = 36k+35$. $36k+35 \equiv 11 \pmod{12}$.
\end{beweis}

\begin{korollar}[Zwingender EA-Schritt nach maximal 2 BC-Schritten]
\label{kor:bca}
Jede Collatz-Bahn, die einen B-Schritt durchläuft, durchläuft nach
spätestens zwei weiteren Schritten einen EA-Schritt. Formal:
Ist $n$ B-Typ, so ist $\U(\U(n))$ stets EA-Typ
(unabhängig davon, ob $n \equiv 7$ oder $n \equiv 19 \pmod{24}$).

\emph{Achtung:} Dies beschreibt nur die mod-24-Ebene. C-Ketten der Länge~$\geq 2$
sind auf feinerer Ebene (mod~$48$, mod~$96$, \ldots) möglich.
\end{korollar}

\begin{beweis}
Aus Prop.~\ref{prop:mod24}: \\
--- Fall B$\to$A ($n \equiv 19 \pmod{24}$): Bereits der erste Folgeschritt ist A-Typ (EA). \\
--- Fall B$\to$C ($n \equiv 7 \pmod{24}$): $n_1 \equiv 11 \pmod{12}$, C-Typ.
Dann gilt $n_1 = 36k+11$ für passendes~$k$.
$n_1 \equiv 11 \pmod{24}$ oder $n_1 \equiv 23 \pmod{24}$.
Im Fall $n_1 \equiv 11 \pmod{24}$ folgt sofort A-Typ.
Im Fall $n_1 \equiv 23 \pmod{24}$ folgt $n_2 \equiv 11 \pmod{12}$ (C-Typ wieder).
\end{beweis}

Die Länge der C-Kette nach dem B-Schritt kann damit größer als~$1$ sein.
Die vollständige Analyse der C-Kettenlängen erfordert die 2-adische Methode.

% ====================================================
\section{Kombinierter Abstiegssatz: BCA-Analyse}
\label{sec:bca}
% ====================================================

\subsection{Exakte Faktorberechnung}

Sei $n_0$ B-Typ mit $n_0 \equiv 7 \pmod{24}$, damit $n_1 = \U(n_0)$ C-Typ.
Sei weiter $n_1 \equiv 11 \pmod{24}$, damit $n_2 = \U(n_1)$ A-Typ.
Ein vollständiger B\textbullet C\textbullet A-Block liefert dann:

\begin{proposition}[Exakte B\textbullet C\textbullet A-Faktoren]
\label{prop:bca-faktoren}
Seien $n_0, n_1, n_2$ wie oben, und sei $\nu := \vii(3n_2+1) \geq 2$.
Dann gilt:
\[
  n_1 = \frac{3n_0 + 1}{2}, \qquad
  n_2 = \frac{9n_0 + 5}{4}, \qquad
  n_3 := \U(n_2) = \frac{27n_0 + 19}{4 \cdot 2^\nu}.
\]
Der Verhältnisfaktor des gesamten Blocks beträgt:
\[
  \frac{n_3}{n_0}
  = \frac{27n_0 + 19}{4 \cdot 2^\nu \cdot n_0}
  \;\approx\; \frac{27}{4 \cdot 2^\nu}
  \quad\text{für großes }n_0.
\]
\end{proposition}

\begin{beweis}
$n_1 = (3n_0+1)/2$ (B-Schritt, $\vii=1$).
$n_2 = (3n_1+1)/2 = (3(3n_0+1)/2+1)/2 = (9n_0+5)/4$ (C-Schritt, $\vii=1$).
$n_3 = (3n_2+1)/2^\nu = (3(9n_0+5)/4+1)/2^\nu = (27n_0+19)/(4\cdot 2^\nu)$.
\end{beweis}

\begin{tcolorbox}[ergebnis]
\textbf{Kontraktionsbedingung für den B\textbullet C\textbullet A-Block:}
\[
  \frac{n_3}{n_0} < 1
  \;\Longleftrightarrow\;
  27n_0 + 19 < 4 \cdot 2^\nu \cdot n_0
  \;\Longleftrightarrow\;
  n_0 > \frac{19}{4 \cdot 2^\nu - 27}.
\]
Das Nennervorzeichen bestimmt die Kontraktionsschwelle:
\begin{itemize}
  \item $\nu = 2$: $4 \cdot 4 = 16 < 27$, \emph{kein Kontraktor möglich}
    (Nenner negativ; $n_3/n_0 \approx 27/16 \approx 1{,}69 > 1$).
  \item $\nu = 3$: $4 \cdot 8 = 32 > 27$, \emph{Kontraktion für} $n_0 > 19/(32-27) = 3{,}8$,
    also für alle $n_0 \geq 5$. Faktor $\approx 27/32 \approx 0{,}84 < 1$.
  \item $\nu = 4$: Faktor $\approx 27/64 \approx 0{,}42 < 1$ (starke Kontraktion).
\end{itemize}
\end{tcolorbox}

\subsection{Wann gilt $\nu \geq 3$ beim A-Schritt?}

Sei $n_2 = (9n_0+5)/4$ A-Typ. Dann $n_2 \equiv 5 \pmod{12}$, und
$3n_2+1 \equiv 16 \equiv 4 \pmod{12}$, also $\vii(3n_2+1) \geq 2$.
Die Frage nach $\nu \geq 3$ hängt davon ab, ob $8 \mid 3n_2+1$:

\begin{proposition}[Bedingung für $\nu \geq 3$ beim A-Schritt]
\label{prop:nu3}
Mit $n_0 \equiv 7 \pmod{24}$ (schreibe $n_0 = 24m+7$) gilt
$n_2 = 54m + 17$ und
\[
  3n_2 + 1 = 162m + 52 = 4(40m + 13) + 4(m\cdot 0 + \tfrac{162m+52-4(40m+13)}{4}).
\]
Präziser: $3n_2 + 1 = 162m + 52$. Hiervon gilt:
\begin{itemize}
  \item $m$ gerade ($m = 2s$): $3n_2+1 = 324s+52 = 4(81s+13)$,
    und $81s+13$ ist ungerade für $s$ gerade, gerade für $s$ ungerade.
    Damit $\vii(3n_2+1) = 2$ für $s$ gerade ($m \equiv 0 \pmod 4$, $n_0 \equiv 7 \pmod{96}$)
    und $\vii \geq 3$ für $s$ ungerade ($n_0 \equiv 55 \pmod{96}$).
  \item $m$ ungerade: entsprechende mod-96-Analyse ergibt analog $\nu \in \{2, 3, \ldots\}$.
\end{itemize}
\end{proposition}

Für die heuristische Abschätzung genügt das empirische Mittel:
Numerisch beobachtet man für A-Typ-Zahlen $n \equiv 5 \pmod{12}$
durchschnittlich $\overline{\vii}(3n+1) \approx 3{,}0$ (vgl.\ \textit{collatz\_eabc\_spektrum.tex}).
Das bedeutet: Im Mittel gilt $\nu = 3$, also \emph{mittlere} Kontraktion
mit Faktor $\approx 27/32 < 1$ pro B\textbullet C\textbullet A-Block.

\subsection{Was nach einem B\textbullet C\textbullet A-Block mit $\nu = 2$ folgt}

Im Fall $\nu = 2$ wächst $n_3 \approx (27/16)n_0$. Der nächste Schritt
startet bei $n_3$, das E-Typ oder A-Typ ist. E-Typ hat typischerweise
$\vii \approx 3$, was starke Kontraktion $\approx 3/8$ liefert. Dann ist
$n_4 \approx (3/8) \cdot n_3 \approx (3/8)(27/16)n_0 = (81/128)n_0 \approx 0{,}63\,n_0 < n_0$.

\begin{tcolorbox}[warnung]
\textbf{Heuristischer Zweischrittabstieg:}
Selbst im ungünstigsten Fall $\nu = 2$ kontrahiert die Bahn
nach einem weiteren EA-Schritt:
\[
  n_0 \;\xrightarrow{\text{B}\cdot\text{C}\cdot\text{A}_2}\;
  \tfrac{27}{16}n_0 \;\xrightarrow{\text{EA}_3}\;
  \tfrac{3}{8}\cdot\tfrac{27}{16}n_0 = \tfrac{81}{128}n_0 < n_0.
\]
Dies ist eine \emph{heuristische} Aussage (mittleres $\vii$), kein Beweis.
Das tatsächliche $\vii$ kann für einzelne $n$ kleiner ausfallen.
\end{tcolorbox}

\subsection{Geometrischer Mittelwert der Schrittfaktoren}

Die numerischen Daten aus \textit{collatz\_eabc\_spektrum.tex} geben:
\[
  \text{B, C:}\; \U(n)/n \approx 3/2,\qquad
  \text{E, A:}\; \U(n)/n \approx 3/8.
\]
Da jede der vier EABC-Klassen gleichgroßen Anteil unter den ungeraden Zahlen hat,
ist der geometrische Mittelfaktor:
\[
  \sqrt[4]{\frac{3}{2} \cdot \frac{3}{2} \cdot \frac{3}{8} \cdot \frac{3}{8}}
  = \sqrt[4]{\frac{9}{4} \cdot \frac{9}{64}}
  = \sqrt[4]{\frac{81}{256}}
  = \frac{3}{4} = 0{,}75 < 1.
\]
Der geometrische Mittelfaktor \emph{aller} odd-to-odd-Schritte
liegt klar unter~$1$. Das erklärt heuristisch die globale Konvergenz
--- ist aber kein Beweis, da einzelne Bahnen vom Mittelwert abweichen können.

% ====================================================
\section{C-Ketten: Endlichkeit und verbleibende Lücke}
\label{sec:cketten}
% ====================================================

\subsection{Die 2-adische Formel}

Das stärkste algebraische Resultat des Beweisversuchs ist die geschlossene
Formel für C-Kettenglieder:

\begin{theorem}[C-Ketten-Formel, aus \textit{collatz\_c\_ketten\_2adisch.tex}]
\label{thm:cformel}
Ist $n_0$ der Startwert einer C-Kette (alle Glieder $n_0,\ldots,n_{k-1}$
sind C-Typ), so gilt für alle $j \geq 0$:
\[
  n_j = \frac{3^j(n_0+1)}{2^j} - 1.
\]
\end{theorem}

Für $n_j \in \Z$ muss $2^j \mid 3^j(n_0+1)$. Da $\gcd(3,2)=1$, folgt:
\[
  2^j \mid (n_0+1) \quad \text{für alle } j \leq k.
\]
Das bedeutet: Eine C-Kette der Länge~$k$ erfordert $\vii(n_0+1) \geq k$.

\subsection{2-adisches Grenzwertargument}

\begin{theorem}[Keine unendlichen C-Ketten in $\N$, 2-adisch]
\label{thm:2adisch}
Eine unendliche C-Kette würde $\vii(n_0+1) \geq k$ für \emph{alle} $k \geq 1$
erfordern, d.\,h.\ $n_0 + 1 \equiv 0 \pmod{2^k}$ für alle $k$.
Das bedeutet: $n_0 = -1$ in den 2-adischen ganzen Zahlen $\Ztwo$.
Da $-1 \notin \N_{\geq 1}$, gibt es in $\N$ kein $n_0$, das eine
unendliche C-Kette startet --- \emph{vorausgesetzt}, dass
der 2-adische Grenzwert $n_0 = -1$ tatsächlich der einzige
Häufungspunkt der Residuenklassenfolge ist.
\end{theorem}

\begin{tcolorbox}[lücke]
\textbf{Verbleibende Lücke beim 2-adischen Argument:}
Theorem~\ref{thm:2adisch} zeigt, dass kein \emph{festes} $n_0 \in \N$
für alle $k$ die Bedingung $n_0 \equiv -1 \pmod{2^k}$ erfüllt.
Das schließt unendliche C-Ketten für jedes einzelne $n_0$ aus.

Was \emph{nicht} direkt folgt: Dass auch C-Ketten beliebig großer
\emph{endlicher} Länge für immer größere $n_0$ tatsächlich
zur Kollatz-Konvergenz keinen Widerspruch erzeugen.
Der kritische Schritt --- von ``keine unendliche C-Kette''
zu ``C-Ketten sind gleichmäßig durch $m_0$ beschränkt'' --- bleibt offen.
\end{tcolorbox}

\begin{tcolorbox}[ergebnis]
\textbf{Was rigoros gilt:}
Für jedes feste $n_0 \in \N$ gibt es nur endlich viele aufeinanderfolgende
C-Schritte (da nur endlich viele $k$ die Bedingung $\vii(n_0+1) \geq k$
erfüllen). Es gibt also keine ``unendliche C-Kette'' im elementarsten Sinn.
Die Frage nach einer \emph{gleichmäßigen} Schranke $m_0$ bleibt jedoch offen.
\end{tcolorbox}

% ====================================================
\section{Zyklenfreiheit: Logarithmus-Argument}
\label{sec:zyklen}
% ====================================================

\subsection{Die Tao-Identität und die Zyklengleichung}

Für die odd-to-odd-Abbildung $\U$ und ein $l$-periodisches Element
$n_0$ (Zyklus $n_0 \to n_1 \to \ldots \to n_{l-1} \to n_0$) lautet
die Tao-Identität:
\[
  n_0 \cdot 2^m = 3^l \cdot n_0 + G_l,
\]
wobei $m = \sum_{i=0}^{l-1} \vii(3n_i+1)$ die Gesamtzahl der Halbierungen ist
und $G_l = \sum_{j=0}^{l-1} 3^{l-1-j} \cdot 2^{\sum_{i=0}^j \vii(3n_i+1)} > 0$.

Da $G_l > 0$ und $n_0 > 0$, folgt sofort die notwendige Bedingung:
\begin{equation}
  \label{eq:zyklus-bed}
  2^m > 3^l,
  \quad\text{d.\,h.}\quad
  \frac{m}{l} > \frac{\log 3}{\log 2} \approx 1{,}58496.
\end{equation}

\subsection{Zyklen mit B\textbullet C\textbullet A-Struktur sind unmöglich}

\begin{theorem}[Keine B\textbullet C\textbullet A-Zyklen]
\label{thm:kein-bca-zyklus}
Sei $l = 3s$ (ein Zyklus, der aus $s$ vollständigen B\textbullet C\textbullet A-Blöcken
besteht, wobei jeder A-Schritt genau $\vii = 2$ Halbierungen liefert).
Dann gilt $m = 4s$ (pro Block: $1 + 1 + 2 = 4$ Halbierungen) und:
\[
  \frac{m}{l} = \frac{4s}{3s} = \frac{4}{3} \approx 1{,}333 < \frac{\log 3}{\log 2} \approx 1{,}585.
\]
Damit ist $2^m = 2^{4s} = 16^s < 27^s = 3^{3s} = 3^l$,
im Widerspruch zur notwendigen Bedingung~\eqref{eq:zyklus-bed}.
\end{theorem}

\begin{beweis}
Jeder der $s$ Blöcke liefert genau $1$ (B) $+$ $1$ (C) $+$ $2$ (A, $\vii=2$)
$= 4$ Halbierungen und $3$ odd-to-odd-Schritte. Also $m = 4s$, $l = 3s$.
$2^{4s} = 16^s$, $3^{3s} = 27^s$, und $16 < 27$ impliziert $16^s < 27^s$. \qed
\end{beweis}

\begin{korollar}[Keine B\textbullet C\textbullet A-Zyklen mit $\vii \in \{2,3\}$]
Für Zyklen mit $s$ B\textbullet C\textbullet A-Blöcken, von denen $s_2$ den Wert $\vii=2$
und $s_3$ den Wert $\vii=3$ beim A-Schritt haben ($s_2+s_3=s$), gilt:
\[
  \frac{m}{l} = \frac{4s_2 + 5s_3}{3s}
  = \frac{4}{3} + \frac{s_3}{3s}
  \leq \frac{4}{3} + \frac{1}{3} = \frac{5}{3} \approx 1{,}667.
\]
Dies übersteigt $\log3/\log2 \approx 1{,}585$ sobald $s_3/s > (1{,}585 - 4/3)/(1/3) = 0{,}753$.
Zyklen sind also ausgeschlossen, wenn weniger als $75{,}3\%$ der A-Schritte
$\vii = 3$ haben.

Wenn alle A-Schritte $\vii = 2$ haben (Minimum), ist $m/l = 4/3 < \log3/\log2$.
Wenn mehr als $75\%$ der A-Schritte $\vii \geq 3$ haben, könnte $m/l > \log3/\log2$
sein --- dann hilft nur noch das Baker-Argument.
\end{korollar}

\subsection{Das Baker-Argument für allgemeine Zyklen}

Für allgemeine Zyklen, bei denen $m/l > \log 3/\log 2$, liefert die
Zyklengleichung:
\[
  n_0 = \frac{G_l}{2^m - 3^l}.
\]
Baker's Theorem über Linearformen in Logarithmen besagt:
\[
  \abs{m \log 2 - l \log 3} > \frac{C}{l^A}
\]
für effektive Konstanten $C, A > 0$. Daraus folgt:
\[
  2^m - 3^l = 3^l\bigl(2^{m-l\log3/\log2} - 1\bigr) \gtrsim 3^l \cdot \frac{C}{l^A}.
\]
Für großes $l$ wächst der Nenner $2^m - 3^l$ exponentiell, während
$G_l < 3^l \cdot n_0 \cdot \text{const}$ (aus der Tao-Identität).
Dies begrenzt $n_0$ nach oben durch einen von~$l$ abhängigen Ausdruck.
In Verbindung mit numerischer Verifikation für kleine Zyklen
(alle $n \leq 2^{68}$ ohne nichtrivialen Zyklus geprüft) liefert das:

\begin{tcolorbox}[konj]
\textbf{Konjektur (bedingt auf Baker-Axiom):}
Für alle $l \geq 2$ und alle $m$ mit $2^m > 3^l$ und $m/l > \log3/\log2$
ist $G_l / (2^m - 3^l)$ keine positive ganze Zahl,
ausgenommen der triviale Zyklus $n_0 = 1$ ($l=1$, $m=2$, $3^1 < 2^2$).

Dies entspricht der konditionalen Zyklenfreiheit in
\texttt{ericmerle3789/collatz-conditional-cycles}.
\end{tcolorbox}

\subsection{Verbindung von Zyklenfreiheit und DC-Struktur}

Der Satz über Zyklenfreiheit hängt eng mit der DC zusammen:

\begin{tcolorbox}[kernidee]
\textbf{Logische Struktur:}
\begin{enumerate}[label=(\arabic*)]
  \item DC (Durchlaufbedingung): Jede Bahn hat beschränkte C-Ketten.
  \item Zyklenfreiheit: Kein nichtrivialer Zyklus.
  \item Keine divergenten Bahnen: Alle Bahnen sind beschränkt.
  \item[(1)+(2)+(3)] $\Rightarrow$ Collatz-Vermutung.
\end{enumerate}
Die DC impliziert, dass $m/l$ nach unten durch ein strukturelles Minimum
begrenzt ist. Da $4/3 < \log3/\log2$ gilt (Theorem~\ref{thm:kein-bca-zyklus}),
schließt die DC die Zyklen mit $m/l \in [4/3,\,\log3/\log2)$ aus.
Für $m/l > \log3/\log2$ braucht man Baker.
\end{tcolorbox}

% ====================================================
\section{Keine divergenten Bahnen}
\label{sec:divergenz}
% ====================================================

\subsection{$E_8$-Normabstieg und Verbindung zur DC}

Aus \textit{collatz\_oktonionen\_beweis.tex} ist bekannt:

\begin{theorem}[Normabstieg für EA-Schritte]
Für $x \in \mathbb{O}_{\Z}$ ungerade mit $\vii(3x+1) \geq 2$ (EA-Schritt)
gilt für $N(x) \gg 1$:
\[
  N(\mathcal{U}_{\mathbb{O}}(x)) < N(x).
\]
Konkret folgt aus $N(3x+1) = 9N(x) + 6\operatorname{Re}(x) + 1$ und
$N(\mathcal{U}_{\mathbb{O}}(x)) = N(3x+1)/2^{2\nu} \leq N(3x+1)/16$,
dass $N(\mathcal{U}_{\mathbb{O}}(x)) \leq (9N(x) + 6\sqrt{N(x)} + 1)/16 < N(x)$
für $N(x) > C_0$ (explizit berechenbar).
\end{theorem}

\begin{tcolorbox}[warnung]
\textbf{Die verbleibende Lücke (Divergenz):}
Aus dem $E_8$-Normabstieg für EA-Schritte folgt: Wenn die DC gilt
(jede Bahn macht nach endlich vielen BC-Schritten einen EA-Schritt),
dann steigt die Norm langfristig nicht. Daraus folgt die Beschränktheit
aller Bahnen für $n > C_0$.

Das Problem ist der Transport der Aussage: Der Normabstieg gilt im
$E_8$-Gitter, und für $n \in \N \subset \R \subset \mathbb{O}_{\Z}$
überträgt sich der Normabstieg auf $n$ (da Einbettungsverträglichkeit bewiesen ist).
Aber: \emph{ohne die DC} weiß man nicht, dass nach endlich vielen Schritten
ein EA-Schritt folgt --- und damit dreht sich das Argument im Kreis.
\end{tcolorbox}

\subsection{Kombiniertes Abstiegsargument (konditional)}

\begin{theorem}[Konditionales Abstiegstheorem]
\label{thm:abstieg}
Vorausgesetzt die DC gilt (jede Bahn macht nach höchstens $m_0$ BC-Schritten
einen EA-Schritt), dann gilt für alle $n > C_0$:
\begin{enumerate}[label=(\roman*)]
  \item Jede Bahn von $n$ ist beschränkt: Es gibt $M(n)$ mit
    $\U^k(n) \leq M(n)$ für alle $k$.
  \item Zusammen mit Zyklenfreiheit (Baker-bedingt) folgt:
    Die Bahn erreicht~$1$.
\end{enumerate}
\end{theorem}

Der Satz ist korrekt, aber seine Voraussetzungen --- DC und Zyklenfreiheit ---
sind selbst die offenen Kernfragen.

% ====================================================
\section{Hauptsatz-Versuch}
\label{sec:hauptsatz}
% ====================================================

Der stärkstmögliche Satz, der aus den bewiesenen Bausteinen folgt,
ohne auf Konjekturen zurückzugreifen, lautet:

\begin{theorem}[Stärkster bewiesener Satz des Beweisversuchs]
\label{thm:haupt}
Sei $n \in \N$ ungerade.
\begin{enumerate}[label=(\arabic*)]
  \item \emph{(BC-Struktursatz)} Auf jeden B-Schritt folgt entweder unmittelbar
    ein EA-Schritt oder genau ein C-Schritt, nach dem ein EA-Schritt folgt
    (auf mod-24-Ebene). Ein zweiter B-Schritt kann auf keinen B-Schritt folgen.

  \item \emph{(C-Ketten-Endlichkeit)} Für jedes feste $n \in \N$ endet
    jede von~$n$ startende C-Kette nach endlich vielen Schritten
    (nämlich nach höchstens $\vii(n+1)$ Schritten).

  \item \emph{(Exponentielle Seltenheit)} Die Menge der Startwerte
    $n \leq X$, von denen C-Ketten der Länge~$\geq k$ starten,
    hat Dichte $\leq 1/2^k$ unter den ungeraden Zahlen.

  \item \emph{(B\textbullet C\textbullet A-Zyklenausschluss)} Zyklen der
    odd-to-odd-Abbildung, die ausschließlich aus B\textbullet C\textbullet A-Blöcken
    mit $\vii = 2$ beim A-Schritt bestehen, existieren nicht.

  \item \emph{(Normabstieg)} Für alle hinreichend großen $n$ gilt:
    Wenn $\vii(3n+1) \geq 2$, so ist $\U(n) < n$.
\end{enumerate}
\end{theorem}

Die Collatz-Vermutung folgt aus diesem Satz \emph{nicht},
da der Übergang von C-Ketten-Endlichkeit zu gleichmäßiger Beschränkung
(DC) nicht gezeigt wurde.

% ====================================================
\section{Ehrliches Fazit}
\label{sec:fazit}
% ====================================================

\subsection{Was tatsächlich bewiesen ist}

\begin{tcolorbox}[ergebnis]
\textbf{Rigoros bewiesen (ohne Axiome außer den Grundlagen der Arithmetik):}
\begin{enumerate}[label=(\arabic*)]
  \item B-Typ kann nicht direkt auf B-Typ folgen
    (Kongruenzrechnung mod~12).
  \item Für jedes feste $n$: die von $n$ startende C-Kette hat Länge
    $\leq \vii(n+1) < \infty$ (2-adisches Divisibilitätsargument).
  \item C-Ketten der Länge~$k$ erfordern Startwerte in Klassen der
    Dichte~$\leq 1/2^k$ (Induktionsbeweis über Residuenbedingungen).
  \item B\textbullet C\textbullet A-Zyklen mit $\vii=2$ beim A-Schritt sind
    ausgeschlossen (Tao-Identität + $16^s < 27^s$).
  \item $E_8$-Normabstieg bei EA-Schritten für große Normen.
\end{enumerate}
\end{tcolorbox}

\subsection{Was offen bleibt --- genau benannt}

\begin{tcolorbox}[lücke]
\textbf{Offene Lücken, exakt benannt:}
\begin{enumerate}[label=(\arabic*)]
  \item \textbf{Gleichmäßige DC-Schranke $m_0$:}
    Es fehlt eine Schranke der Form ``jede Bahn macht nach höchstens $m_0$
    Schritten einen EA-Schritt für ein \emph{von $n$ unabhängiges} $m_0$''.
    Hier gibt das 2-adische Argument nur eine $n$-abhängige Schranke.

  \item \textbf{Zyklenfreiheit für alle Zyklen:}
    Zyklen mit $m/l > \log3/\log2$ (hoher EA-Anteil) sind nur bedingt
    (Baker-Axiom) ausgeschlossen. Ohne Baker ist Zyklenfreiheit unbewiesen.

  \item \textbf{Verbindung 2-adisch $\to$ universell:}
    Das 2-adische Argument zeigt: kein $n_0$ kann unendlich lang eine
    C-Kette bilden. Es zeigt \emph{nicht}: C-Ketten sind gleichmäßig beschränkt.
    Der Unterschied: existenzielle vs.\ universelle Aussage über Kettenlängen.

  \item \textbf{Keine divergenten Bahnen:}
    Der $E_8$-Normabstieg bei EA-Schritten setzt die DC voraus und ist
    daher zirkulär, bis die DC eigenständig bewiesen ist.

  \item \textbf{Kleinstfall:}
    Die strukturellen Argumente gelten für ``hinreichend große'' $n$.
    Für kleine $n$ ($n \leq C_0$) ist numerische Verifikation nötig,
    die für $n \leq 10^7$ (dieses Projekt) bzw.\ $n \leq 2^{68}$
    (Stand der Literatur) vorliegt.
\end{enumerate}
\end{tcolorbox}

\subsection{Gesamteinschätzung}

\begin{tcolorbox}[kernidee]
\textbf{Wo steht dieser Beweisversuch?}

Die fünf Vorgängerdokumente haben eine kohärente algebraisch-modulare
Infrastruktur aufgebaut, die reale und rigorose Ergebnisse enthält.
Die stärksten davon --- B-Ketten-Ausschluss, C-Ketten-Formel,
$m/l$-Argument gegen B\textbullet C\textbullet A-Zyklen --- sind
korrekte Lemmata, die auch in formalen Beweissystemen verifikationsfähig sind.

Die \emph{eigentliche} Collatz-Lücke, nämlich die DC, ist durch diese
Methoden nicht schließbar. Die DC ist äquivalent zu einer gleichmäßigen
Schranke über alle Orbits zugleich --- genau die Art von globaler
Aussage, die sich nicht aus lokalen modularen Beobachtungen gewinnen lässt.

Der erfolgversprechendste nächste Schritt bleibt:
\begin{itemize}
  \item Ergodentheoretische Methoden (Taos Martingal-Ansatz),
    spezialisiert auf C-Ketten.
  \item Transferoperator-Spektralanalyse für die C-Ketten-Unterabbildung.
  \item Baker-Schranken für die Zyklengleichung, formalisiert in Lean\,4.
\end{itemize}

\textbf{Klare Schlussfolgerung:} Dieser Beweisversuch liefert keine
Lösung der Collatz-Vermutung, aber eine saubere, rigorose Reduktion
auf die Frage: ``Kann eine Collatz-Bahn beliebig lange C-Ketten
\emph{gleichmäßig} bilden?'' Diese Reduktion ist selbst ein Ergebnis.
\end{tcolorbox}

\subsection{Statusmatrix: Gesamtübersicht}

\begin{center}
\renewcommand{\arraystretch}{1.5}
\begin{tabular}{p{6.5cm} p{2.8cm} p{4cm}}
\toprule
\textbf{Aussage} & \textbf{Status} & \textbf{Bemerkung} \\
\midrule
B$\to$B unmöglich & \textcolor{bewiesengrün}{\bfseries Bewiesen} & mod 12 \\
B$\to$C oder B$\to$A (mod 24) & \textcolor{bewiesengrün}{\bfseries Bewiesen} & mod 24 \\
C-Ketten-Formel & \textcolor{bewiesengrün}{\bfseries Bewiesen} & Induktion \\
C-Kette Länge $\leq \vii(n_0+1)$ & \textcolor{bewiesengrün}{\bfseries Bewiesen} & 2-adisch \\
Exponentielle C-Ketten-Seltenheit & \textcolor{bewiesengrün}{\bfseries Bewiesen} & Residuenklassen \\
B\textbullet C\textbullet A-Zyklus ausgeschlossen & \textcolor{bewiesengrün}{\bfseries Bewiesen} & Tao-Ident. + $16<27$ \\
$E_8$-Normabstieg (EA-Schritte) & \textcolor{bewiesengrün}{\bfseries Bewiesen} & groß $n$, mod bedingt \\
DC (gleichmäßige Schranke $m_0$) & \textcolor{lückenrot}{\bfseries Offen} & Kernlücke \\
Zyklenfreiheit allgemein & \textcolor{warnorange}{\bfseries Bedingt} & Baker-Axiom \\
Keine divergenten Bahnen & \textcolor{warnorange}{\bfseries Bedingt} & setzt DC voraus \\
Collatz-Vermutung & \textcolor{lückenrot}{\bfseries Offen} & --- \\
\bottomrule
\end{tabular}
\end{center}

% ====================================================
% Literatur
% ====================================================
\begin{thebibliography}{9}

\bibitem{collatz-eabc}
T. Hoffbauer,
\textit{Die verborgene Binnenstruktur der Collatz-Dynamik:
Eine Spektrum-Notiz zur odd-to-odd-Abbildung im EABC-Rahmen},
Mathematische Texte (2025).

\bibitem{collatz-oktonionen}
T. Hoffbauer,
\textit{Collatz-Dynamik auf ganzzahligen Oktonionen:
Ein Beweisversuch über das $E_8$-Gitter},
Mathematische Texte (2025).

\bibitem{collatz-dc}
T. Hoffbauer,
\textit{Die Durchlaufbedingung im Morley-Walter-Rahmen},
Mathematische Texte (2025).

\bibitem{collatz-c-ketten}
T. Hoffbauer,
\textit{C-Ketten in der Collatz-Dynamik:
Eine 2-adische Analyse der Durchlaufbedingung},
Mathematische Texte (2025).

\bibitem{collatz-lean}
T. Hoffbauer,
\textit{Lean\,4-Formalisierung der Durchlaufbedingung},
Mathematische Texte (2026).

\bibitem{tao-collatz}
T. Tao,
\textit{Almost all orbits of the Collatz map attain almost bounded values},
Forum of Mathematics Pi \textbf{10} (2022), e12.

\bibitem{baker}
A. Baker,
\textit{Linear forms in the logarithms of algebraic numbers},
Mathematika \textbf{13} (1966), 204--216.

\bibitem{lagarias}
J.\,C. Lagarias,
\textit{The $3x+1$ problem and its generalizations},
Amer.\ Math.\ Monthly \textbf{92} (1985), 3--23.

\end{thebibliography}

\end{document}
